\newglossaryentry{glos:elo}
{
    name=ELO,
    description={Eine Vergleichszahl, um die Stärke von Spieler untereinander zu vergleichen (\cite{elo})}
}
\newglossaryentry{glos:fide}
{
    name=FIDE,
    description={Die Fédération Internationale des Échecs (FIDE) ist der Dachverband des Schachsports und regelt alle internationalen Schachwettbewerbe. Als nichtstaatliche Institution gegründet, wurde sie 1999 vom Internationalen Olympischen Komitee als globale Sportorganisation anerkannt (\cite{fideAbout})}
}
\newglossaryentry{glos:ccrl}
{
    name=CCRL,
    description={Die Computer Chess Rating Lists (CCRL) ist eine in 2005 gegründete Gruppe, welche Schachcomputer in unterschiedlichen Modi gegeneinander spielen lassen (\cite{ccrlAbout})}
}
\newglossaryentry{glos:cegt}
{
    name=CEGT,
    description={Die Chess Engines Grand Tournament (CEGT) ist eine in 2001 gegründete Gruppe. Sie vergleichen Schachcomputer in unterschiedlichen Modi (\cite{cegtWebseite})}
}
\newglossaryentry{glos:tcec}
{
    name=TCEC,
    description={Die Top Chess Engine Championship (TCEC) ist ein Computerschach-Turnier, das von Chessdom in Zusammenarbeit mit der Chessdom Arena organisiert wird. Besonders ist, dass die Spiele online übertragen werden und die Computer eine Bedenkzeit haben (\cite{tcecWebseite})}
}
\newglossaryentry{glos:stockfish}
{
    name=Stockfish,
    description={Stockfish ist eine Open‑Source‑Schachengine (\cite{stockfishGithub})}
}
\newglossaryentry{glos:JUnit}
{
    name=JUnit,
    description={JUnit ist ein Java‑Testing‑Framework (\cite{junitHomepage})}
}